\Theta'_{17}=1,\quad
\Theta'_{19}=37,\quad
\Theta'_{23}=47^2\cdot137.
\]
For comparison $\Xi'_{23}=47\cdot137$: split prime support agrees, valuations
differ.  The reference SymPy implementation reduces the $K=17$ computation
from roughly $51$ seconds for the full shift resultant to roughly $0.6$ seconds
for the translated half-degree route; at $K=23$ the latter remains on the
order of seconds, after which Smith normal form becomes the principal finite
algebra bottleneck.  These timings are reproducibility data, not asymptotic
complexity theorems.

The norm of the inverse constant shift coefficient also has the exact shape
\begin{equation}\label{eq:inverse-norm}
\left|\Norm(R_-(0))\right|
 =2^{m^2}K^m\left|\Norm(\Res(E,O))\right|^2,
\end{equation}
using $|\Norm(F(0))|=K^m$.  Hence every prime away from $2K$ occurs with even
valuation in this particular norm.  This explains part of the observed square
pattern but does not imply equality of the support integers.

\section{Global split-prime budget and minimal-extension Frobenius}

Simultaneous multiplication of collision exponent indices by
$a\in\mathbf F_K^\times$ acts freely on non-forced collision quadruples.  The
number of normalized scaling orbits is
\[
N_K=K^3-5K^2+6K+1.
\]
Combining the square-resultant theorem with $|\Res(\Phi_K,P)|\le12^{K-1}$
gives:

\begin{theorem}[Weighted split-prime budget]\label{thm:budget}
\begin{equation}\label{eq:budget}
\sum_{p\equiv1\,(K)}\Delta_K(p)\log p
\le N_K\bigl((K-1)\log12-3\log K\bigr).
\end{equation}
\end{theorem}

For prime $K\asymp X^\theta$ and split prime $p\asymp X$,
Bombieri--Vinogradov and PNT \cite{Bombieri1965,IK2004} imply
\[
\operatorname{average}\frac{\Delta_K(p)}{K^2}
 \ll_\theta X^{3\theta-1}
\]
for fixed $\theta<1/3$.  Integrality of the excess strengthens this to a
density-one exact baseline for fixed $\theta<1/5$:
\[
\Delta_K(p)=0,
\qquad
E_+(A)=3(K-1)(K-2)
\]
for asymptotically almost every split pair.

Now let $p\ne2,K$ be arbitrary and put
\[
f=\ord_K(p),\qquad q=p^f.
\]
In $\mathbf F_q$ the order-$K$ subgroup exists uniquely.

\begin{theorem}[Frobenius amplification]\label{thm:frob}
Every non-forced collision orbit under the $p$-power Frobenius has size exactly
$f$.  Hence
\[
f\mid\Delta_K^{\rm ext}(p).
\]
The weighted resultant budget and the discriminant classification extend to
all such rational primes in the minimal extension.
\end{theorem}

\begin{proof}
Frobenius multiplies every exponent by $p$ modulo $K$.  If a positive power
$d<f$ fixed a collision tuple, any nonzero coordinate would imply
$p^d\equiv1\pmod K$, contradicting the definition of $f$.  The resultant
argument is applied at one prime ideal over $p$; $p\ne K$ is unramified, so its
valuation on rational integers is the ordinary $p$-valuation.
\end{proof}

PNT, Brun--Titchmarsh and the scarcity of small multiplicative orders then give
a density-one exact baseline for unrestricted prime pairs
$K\asymp X^\theta$, $p\asymp X$ for every fixed $\theta<1/3$.  This mechanism
is absent on the split slice because there $f=1$.

\section{Full Cayley--Jacobi trace and corrected off/off geometry}

Let $S=H^\perp$ be the multiplicative-character annihilator and define the
normalized twisted-Jacobi matrix
\[
M_S(q)=\frac{e_p(-q/2)}{p-1}J_S(q).
\]
Character orthogonality gives the exact trace identity
\begin{equation}\label{eq:jacobi}
\widehat{1_A}(r)=K\operatorname{Tr}M_S(2r),
\end{equation}
and therefore
\begin{align}
\sum_q|\operatorname{Tr}M_S(q)|^2&=\frac{p(K-1)}{K^2},\label{eq:trace2}\\
\sum_q|\operatorname{Tr}M_S(q)|^4&=\frac{pE_+(A)}{K^4}.\label{eq:trace4}
\end{align}
